(1) 理想迈克耳孙干涉仪两臂在同一无耗散均匀介质（折射率 $n$）中，输入单色光真空波长 $\lambda_0$。一臂反射镜沿光轴移动 $\Delta L$，求往返光程差变化。
(2) 推导两束光的相位差变化与记录到的完整条纹数。
(3) 镜子缓慢匀速移动速度 $v$ 时，求探测器的条纹扫过频率；对 $n=1$ 与反射多普勒效应作比较。
(4) 若光源不是完全单色，说明为何在光程差远大于相干长度时不能简单按相位公式保持高对比条纹。
